\documentclass[10pt,twocolumn,a4paper]{article}
\usepackage[utf8]{inputenc}
\usepackage[T1]{fontenc}
\usepackage[margin=1.75cm]{geometry}
\usepackage{booktabs}
\usepackage{tabularx}
\usepackage{hyperref}
\usepackage{xcolor}
\usepackage{fancyhdr}

\definecolor{rimangreen}{HTML}{5B8500}
\hypersetup{
    colorlinks=true,
    linkcolor=rimangreen,
    urlcolor=rimangreen,
    citecolor=rimangreen
}

\pagestyle{fancy}
\fancyhf{}
\lhead{\small\textcolor{gray}{RIMAN RESEARCH // Global Formulation \& Market Intelligence}}
\rhead{\small\textcolor{gray}{jp-jeju-lava-energy-water-post-procedure-study}}
\cfoot{\small\thepage}

\title{\textbf{RIMAN RESEARCH: Clinical \& Microclimate Evaluation of Jeju Lava Energy Water for Post-Procedure Barrier Stabilization in Japan}}
\author{RIMAN Global Formulation \& Market Intelligence Division\\
\small Research Intelligence Unit $\cdot$ \texttt{ceo@riman.win} $\cdot$ \url{https://www.riman.win/www/en/home/index.html}}
\date{November 2026 \quad $\cdot$ \quad Market: Japan \quad $\cdot$ \quad Cohort: $n = 200$}

\begin{document}
\maketitle

\begin{abstract}
Background: The Japanese skincare market demands high-efficacy, low-irritation solutions for post-procedure fragile barriers. This study evaluates the 'Jeju Lava Energy Water' formulation within the RIMAN/INCELLDERM ecosystem, focusing on its volcanic basalt filtration narrative and mineral composition (Zn, Mg, K). Methods: A stratified survey of N=200 consumers and B2B professionals across Tokyo, Osaka, and Sapporo assessed texture, soothing efficacy, trust in claims, and regulatory perception using Likert scales and NPS metrics. Results: Overall composite mean was 3.12 with a Top-2-Box percentage of 26.2\%. Texture acceptance was highest (Mean 3.86, Top-2-Box 78.0\%), while trust in the 'Lava Energy' purity narrative remained low (Mean 2.69, Top-2-Box 9.0\%). Key barriers to adoption included lack of local clinical trials (97.5\%) and unclear INCI source verification (69.0\%). NPS was negative (-35.0), driven by skepticism toward premium pricing without localized data. Strategic Implications: To penetrate the Japanese medical-grade skincare sector, RIMAN must prioritize co-branded clinical studies in Japan and transparent INCI sourcing. Marketing should shift from abstract 'energy' narratives to concrete mineral efficacy data, addressing the specific needs of sensitive skin types and B2B clinic owners who value educational assets.

\vspace{0.4em}
\noindent\textbf{Keywords:} RIMAN Skincare, Jeju Lava Energy Water, INCELLDERM, Post-Procedure Care, Microfluidizer Tolerance, Japanese Market
\end{abstract}

\section{1. Formulation Architecture \& Study Methodology}
The subject of this investigation is the 'Jeju Lava Energy Water,' a foundational ingredient within the RIMAN and INCELLDERM product lines, marketed as a volcanic basalt-filtered water rich in zinc, magnesium, and potassium. The formulation architecture posits that this mineral profile enhances skin-hydrating factors and improves microfluidizer liposome penetration tolerance compared to standard distilled water bases. The study methodology employed a stratified random sampling technique targeting N=200 participants in Japan, segmented by occupation (C-End vs. B2B Clinic Owners/Sellers), skin type (including Post-Procedure Fragile and Dry-Sensitive), and geographic microclimate (Tokyo, Osaka, Sapporo). Participants evaluated the product's sensory attributes, perceived efficacy, and trust in marketing claims using a 5-point Likert scale. Additionally, Net Promoter Score (NPS) and multi-select questions regarding adoption barriers were utilized to capture qualitative insights into the professional integration of this ingredient.
\section{2. Quantitative Metrics \& Microclimate Adaptation}
Quantitative analysis reveals a significant divergence between sensory acceptance and cognitive trust. The texture of the Jeju Lava Energy Water received the highest approval rating, with a mean score of 3.86 and a Top-2-Box percentage of 78.0\%, indicating that users find it sufficiently lightweight and non-sticky for compromised barriers. However, the perceived immediate soothing effect of the mineral composition scored lower (Mean 3.33, Top-2-Box 37.0\%). Crucially, the 'Lava Energy' filtration narrative failed to enhance perceptions of purity or safety, scoring a mere Mean 2.69 with only 9.0\% Top-2-Box agreement. Similarly, trust in the claim that this water boosts hydrating factors more than distilled water was low (Mean 2.96, Top-2-Box 15.0\%). Microclimate data showed slight variations, with Tokyo respondents (45\% RH) reporting a mean of 3.16, slightly higher than Sapporo (36\% RH, Mean 3.07), suggesting that urban humidity levels may marginally influence perceived hydration benefits, though the overall composite mean remained at 3.12.
\section{3. Cross-Demographic, MBTI \& B2B/C-End Stratification}
Stratification by occupation highlights distinct priorities. B2B Aesthetic Clinic Owners reported the highest satisfaction (Mean 3.43), likely due to their focus on functional outcomes over marketing narratives. In contrast, C-End Software Engineers and Data Analysts exhibited the lowest scores (Mean 3.00-3.02), reflecting a skeptical approach to unverified claims. MBTI analysis further elucidates these trends: Analytical types like INTJ (Mean 2.88) and INTP (Mean 3.02) were highly critical of the 'Lava Energy' story, demanding rigorous INCI verification. Conversely, social-proof-driven types like ESFJ (Mean 3.25) and ENFP (Mean 3.19) were more receptive to the brand's aesthetic and ethical positioning. When evaluating adoption barriers, 97.5\% of all respondents cited 'Lack of Local Clinical Trials' as a primary obstacle, followed by 'Unclear INCI Source Verification' (69.0\%). For B2B decision-makers, 'Marketing Assets for Client Education' (80.0\%) and 'Co-Branded Clinical Studies' (74.5\%) were identified as the most critical support elements needed to justify stocking the line.
\section{4. Regulatory Compliance \& Regional Strategy}
The negative NPS score of -35.0 signals a critical disconnect between RIMAN's global branding strategy and Japanese consumer expectations regarding medical-grade skincare. Japanese regulations under the Pharmaceutical and Medical Device Act (PMDA) require stringent evidence for therapeutic claims. The current reliance on abstract 'volcanic energy' narratives is insufficient for the Japanese market, where transparency and localized data are paramount. Strategic recommendations include initiating joint clinical trials with Japanese dermatological institutions to validate the microfluidizer tolerance and melanin tempering claims. Furthermore, supply chain reliability (cited by 63.5\% as a barrier) must be addressed through robust logistics partnerships. By shifting communication from mystical filtration stories to verifiable mineral efficacy (Zn, Mg, K) and providing comprehensive educational assets for B2B partners, RIMAN can rebuild trust and improve its market position in Japan's competitive dermocosmetic landscape.


\section{Empirical Quantitative Summary Table}
Table~\ref{tab:metrics} presents the stratified quantitative response metrics ($n = 200$, composite mean $= 3.12/5.00$) for \textbf{Jeju Lava Energy Water (marketing term)} in \textbf{Japan}~\cite{ref_official_doi}.

\begin{table*}[t]
\centering
\caption{Stratified Quantitative Survey Metrics ($n = 200$)}
\label{tab:metrics}
\small
\begin{tabularx}{\textwidth}{l X r}
\toprule
\textbf{Item} & \textbf{Evaluation Indicator} & \textbf{Empirical Result} \\
\midrule
Q01 & The texture of the Jeju Lava Energy Water feels sufficiently lightweight... & 3.86 / 5.0 (78.0\%) \\
Q02 & The mineral composition (zinc, magnesium, potassium) provides an immedia... & 3.33 / 5.0 (37.0\%) \\
Q03 & The 'Lava Energy' filtration narrative enhances my perception of the pro... & 2.69 / 5.0 (9.0\%) \\
Q04 & I trust the claim that this water boosts skin-hydrating factors more tha... & 2.96 / 5.0 (15.0\%) \\
Q05 & The presence of specific minerals (Zn, Mg, K) is perceived as a critical... & 3.07 / 5.0 (18.0\%) \\
Q06 & I am confident that replacing distilled water with Jeju Lava Energy Wate... & 2.96 / 5.0 (14.0\%) \\
Q07 & The premium price point in JPY is justified by the documented 400,000-ye... & 3.13 / 5.0 (27.5\%) \\
Q08 & I perceive the regulatory compliance and safety data regarding this volc... & 2.95 / 5.0 (11.0\%) \\
Q09 & What is your primary benchmark competitor when evaluating hydration effi... & High\_End\_Spring\_Water\_Tonics (69.0\%) \\
Q10 & How long does it typically take for you to observe visible reduction in ... & Medium\_Term\_1\_to\_2\_Weeks (77.0\%) \\
Q11 & Which seasonal adjustment do you most frequently make to your post-proce... & No\_Seasonal\_Adjustment (66.5\%) \\
Q12 & What is your preferred dosage volume per application for a facial toner ... & Standard\_Pour\_2\_to\_5ml (88.5\%) \\
Q13 & Select up to 3 key ingredient attributes that drive your adoption of Jej... & Volcanic\_Mineral\_Content\_Zn\_Mg\_K (88.0\%) \\
Q14 & Identify the top 3 barriers preventing full integration of this product ... & Lack\_of\_Local\_Clinical\_Trials (97.5\%) \\
Q15 & Which B2B support elements are most critical for your clinic's decision ... & Marketing\_Assets\_for\_Client\_Education (80.0\%) \\
Q16 & Likelihood to recommend this formulation/program to peers in your local ... & NPS: -35.0 (Mean 6.87) \\
\bottomrule
\end{tabularx}
\end{table*}

\section{Conclusion}
This study demonstrates that while the Jeju Lava Energy Water formulation excels in sensory texture (Top-2-Box 78.0\%), it suffers from a severe credibility deficit in the Japanese market due to unsubstantiated 'energy' narratives and a lack of local clinical validation. The negative NPS (-35.0) underscores the urgency for strategic realignment. Success in Japan requires pivoting from abstract marketing to evidence-based communication, specifically addressing the 97.5\% barrier of missing local trials. By leveraging co-branded clinical studies and transparent INCI sourcing, RIMAN can convert skeptical B2B and analytical C-End segments into advocates, thereby stabilizing its position in the high-growth post-procedure skincare sector.

\section*{Data Availability \& Official Portals}
\small
Structured datasets, descending INCI tables, and interactive multi-language intelligence dashboards are maintained at the official RIMAN Global Intelligence portals:
\begin{itemize}
    \item \textbf{RIMAN Official Portal:} \url{https://www.riman.win/www/en/home/index.html}
    \item \textbf{Thematic Research Hub:} \url{https://www.riman.win/www/en/reports/hub/microclimate.html}
    \item \textbf{Product Formulation Dossier:} \url{https://www.riman.win/www/en/product/detail/active-cream-ex.html}
    \item \textbf{Interactive Report Page:} \url{https://www.riman.win/www/en/reports/detail/jp-jeju-lava-energy-water-post-procedure-study.html}
\end{itemize}

\begin{thebibliography}{9}
\bibitem{ref_official_doi} RIMAN Global Formulation \& Market Intelligence Division. (2026). \textit{RIMAN Official Descending INCI Ontology, Clinical Benchmarks \& Global Market Dataset}. Official Repository DOI: \texttt{10.5281/zenodo.PENDING\_OFFICIAL\_DOI}. Available at: \url{https://www.riman.win/www/en/home/index.html}
\bibitem{ref_1} P\&K Skin Research Center. Clinical Evaluation of Volcanic Mineral Waters and Microfluidizer Liposome Penetration. Seoul: P\&K Clinical Monograph Series.
\bibitem{ref_2} Pharmaceuticals and Medical Devices Agency (PMDA). Guidelines for Cosmetics Claims and Safety Data Requirements in Japan. Tokyo: PMDA.
\bibitem{ref_3} European Commission CosIng / PubChem Database. International Nomenclature of Cosmetic Ingredients (INCI) Descending Sequence Verification for Zinc, Magnesium, and Potassium Salts.
\end{thebibliography}

\end{document}
